\documentclass[11pt]{article}
\usepackage{reusv}
\usepackage{listings}
\usepackage{enumitem}

\lstset{basicstyle=\ttfamily\small, frame=single, breaklines=true}
\setborderthickness{0.5pt}
\setbordermargin{1cm}
\setbordercolor{black}

\slimtitle{C2. 초기값 생성 전략}

\begin{document}
\drawborder
\thispagestyle{firstpage}

\lightertitle{C2. 초기값 생성 전략}
\def\lighterdate{February 12, 2026}
\lighterauthor{Suwon Lee}
\lighteraddress{}{Future Mobility Control Lab, Department of Future Mobility, Kookmin University, Seoul, Republic of Korea}
\lighteremail{suwon@kookmin.ac.kr}

\begin{abstract}
연속 추력 궤적 최적화에서 NLP 솔버의 수렴 성공 여부는 초기값(initial guess)의 품질에 결정적으로 의존한다.
본 보고서는 두 가지 초기값 생성 전략---선형 보간법(cosine blending 기반)과 케플러 전파법---의 수학적 배경,
구현 매핑, 그리고 수치 검증 결과를 기술한다.
Universal Variable Formulation에 의한 케플러 전파, cosine blending에 의한 출발/도착 궤적 혼합,
그리고 true anomaly에 대한 물리적 휴리스틱을 포함한다.
\end{abstract}

% 대응 코드: \texttt{src/orbit\_transfer/optimizer/initial\_guess.py}

%% =========================================================================
\section{목적}
\label{sec:purpose}
%% =========================================================================

연속 추력 궤적 최적화는 비선형 프로그래밍(NLP)으로 정식화되며, Interior-Point 또는 SQP 계열 솔버로 풀린다.
이러한 솔버는 Newton-Raphson 기반 반복법이므로, 초기값(initial guess)의 품질이 수렴 성공 여부와 계산 효율에 결정적 영향을 미친다.

\textbf{수렴 반경 문제.}
Newton-Raphson 반복의 국소 수렴 이론에 따르면, 해 $\mathbf{x}^*$ 주변의 수렴 반경(convergence radius) $\delta > 0$가 존재하여,
$\|\mathbf{x}_0 - \mathbf{x}^*\| < \delta$인 초기값 $\mathbf{x}_0$에서만 이차 수렴이 보장된다~\cite{Nocedal2006}.
궤적 최적화에서 결정 변수의 차원은 수백\textasciitilde수천에 달하므로, 무작위 초기값은 수렴 반경 밖에 놓일 가능성이 높다.

\textbf{물리적 타당성.}
초기값이 물리법칙에 부합하는 궤적에 가까울수록, NLP의 동역학 구속조건 위반(constraint violation)이 작아져
솔버의 실행가능성 복원(feasibility restoration) 단계가 줄어든다~\cite{Wachter2006}.
케플러 궤도를 기반으로 한 초기값은 무추력 상태에서의 물리적 해이므로, 저추력 문제에서 좋은 출발점을 제공한다.

본 모듈은 두 가지 초기값 생성 전략을 구현한다.
\begin{enumerate}[nosep]
  \item \textbf{선형 보간법}(linear interpolation guess): 출발/도착 궤도를 각각 케플러 전파한 후 cosine blending으로 혼합
  \item \textbf{케플러 전파법}(Keplerian guess): 출발 궤도에서 단순 케플러 전파
\end{enumerate}
두 방법 모두 true anomaly 초기값에 대한 물리적 휴리스틱을 포함한다.


%% =========================================================================
\section{수학적 배경}
\label{sec:math}
%% =========================================================================

%% -------------------------------------------------------------------------
\subsection{케플러 전파: Universal Variable Formulation}
\label{sec:kepler}
%% -------------------------------------------------------------------------

케플러 전파(Kepler propagation)는 이체 문제에서 초기 상태벡터 $(\mathbf{r}_0, \mathbf{v}_0)$로부터 시간 $\Delta t$ 후의 상태 $(\mathbf{r}, \mathbf{v})$를 계산하는 과정이다.
Universal Variable Formulation은 타원, 쌍곡선, 포물선 궤도를 단일 알고리즘으로 처리한다~\cite{Curtis2020}.

\textbf{Stumpff 함수.}
Universal variable $\chi$와 $\psi = \alpha \chi^2$ ($\alpha = 1/a$, 장반경의 역수)에 대해 Stumpff 함수 $c_2(\psi)$, $c_3(\psi)$를 정의한다.
\[
c_2(\psi) = \begin{cases}
\dfrac{1 - \cos\sqrt{\psi}}{\psi} & \psi > 0 \text{ (타원)} \\[6pt]
\dfrac{\cosh\sqrt{-\psi} - 1}{-\psi} & \psi < 0 \text{ (쌍곡선)} \\[6pt]
\dfrac{1}{2} & \psi = 0 \text{ (포물선)}
\end{cases}
\]
\[
c_3(\psi) = \begin{cases}
\dfrac{\sqrt{\psi} - \sin\sqrt{\psi}}{\psi\sqrt{\psi}} & \psi > 0 \\[6pt]
\dfrac{\sinh\sqrt{-\psi} - \sqrt{-\psi}}{(-\psi)\sqrt{-\psi}} & \psi < 0 \\[6pt]
\dfrac{1}{6} & \psi = 0
\end{cases}
\]
이들은 $\cos$ 및 $\cosh$의 Taylor 전개에서 자연스럽게 나타나는 함수이며, $\psi \to 0$에서 $c_2 \to 1/2$, $c_3 \to 1/6$으로 연속이다.

\textbf{역수 장반경.}
초기 상태벡터로부터 역수 장반경 $\alpha$를 계산한다.
\[
\alpha = \frac{2}{r_0} - \frac{v_0^2}{\mu}
\]
$\alpha > 0$이면 타원, $\alpha < 0$이면 쌍곡선, $\alpha = 0$이면 포물선 궤도이다.
본 프로젝트의 원궤도($e = 0$)에서는 항상 $\alpha > 0$이다.

\textbf{Universal Kepler 방정식.}
전파 시간 $\Delta t$에 대응하는 universal variable $\chi$를 다음 방정식의 근으로 구한다.
\begin{equation}
F(\chi) = \chi^3 c_3(\psi) + \frac{\mathbf{r}_0 \cdot \mathbf{v}_0}{\sqrt{\mu}} \chi^2 c_2(\psi) + r_0 \chi (1 - \psi\, c_3(\psi)) - \sqrt{\mu}\,\Delta t = 0
\label{eq:universal_kepler}
\end{equation}
여기서 $\psi = \alpha \chi^2$이고, $r_0 = \|\mathbf{r}_0\|$이다.

\textbf{Newton-Raphson 반복.}
$F(\chi)$의 도함수는 다음과 같이 주어진다.
\[
F'(\chi) = \chi^2 c_2(\psi) + \frac{\mathbf{r}_0 \cdot \mathbf{v}_0}{\sqrt{\mu}} \chi (1 - \psi\, c_3(\psi)) + r_0 (1 - \psi\, c_2(\psi))
\]
이 값은 전파 후 거리 $r(\chi)/\sqrt{\mu}$와 같다. 반복은 $|\delta\chi| < \varepsilon$ ($\varepsilon = 10^{-12}$)일 때 종료한다.
타원 궤도의 초기 추정값은 $\chi_0 = \sqrt{\mu}\,\Delta t\,\alpha$로 설정한다.

\textbf{Lagrange 계수.}
수렴한 $\chi$로부터 Lagrange 계수를 계산한다.
\[
f = 1 - \frac{\chi^2}{r_0} c_2(\psi), \qquad
g = \Delta t - \frac{\chi^3}{\sqrt{\mu}} c_3(\psi)
\]
\[
\dot{f} = \frac{\sqrt{\mu}}{r\,r_0} \chi (\psi\, c_3(\psi) - 1), \qquad
\dot{g} = 1 - \frac{\chi^2}{r} c_2(\psi)
\]
여기서 $r = \|\mathbf{r}\|$이다. 최종 상태벡터는 다음과 같다.
\[
\mathbf{r} = f\,\mathbf{r}_0 + g\,\mathbf{v}_0, \qquad
\mathbf{v} = \dot{f}\,\mathbf{r}_0 + \dot{g}\,\mathbf{v}_0
\]
Lagrange 계수는 항등식 $f\dot{g} - \dot{f}g = 1$을 만족하며,
이는 이체 문제에서 상태 전이의 정준적(canonical) 성질에 기인한다~\cite{Curtis2020}.


%% -------------------------------------------------------------------------
\subsection{선형 보간법 (Cosine Blending)}
\label{sec:cosine_blending}
%% -------------------------------------------------------------------------

출발 궤도와 도착 궤도 각각을 케플러 전파하여 두 개의 참조 궤적을 생성한 후, 이를 시간에 따라 부드럽게 혼합한다.

\textbf{출발 궤도 전파 (forward).}
출발 궤도요소 $(a_0, e_0, i_0, \Omega_0, \omega_0, \nu_0)$에서 초기 상태벡터 $(\mathbf{r}_0, \mathbf{v}_0)$를 계산하고,
각 시간점 $t_k$ ($k = 0, 1, \ldots, N-1$)에서 케플러 전파한다.
\[
(\mathbf{r}_k^{\text{dep}}, \mathbf{v}_k^{\text{dep}}) = \text{Kepler}(\mathbf{r}_0, \mathbf{v}_0, t_k)
\]

\textbf{도착 궤도 역전파 (backward).}
도착 궤도요소 $(a_f, e_f, i_f, \Omega_f, \omega_f, \nu_f)$에서 최종 상태벡터 $(\mathbf{r}_f, \mathbf{v}_f)$를 계산하고,
각 시간점에서 역방향으로 전파한다. 역전파는 $\Delta t = t_k - T < 0$을 전파 시간으로 사용하여 수행한다.
\[
(\mathbf{r}_k^{\text{arr}}, \mathbf{v}_k^{\text{arr}}) = \text{Kepler}(\mathbf{r}_f, \mathbf{v}_f, t_k - T)
\]

\textbf{Cosine blending.}
혼합 가중치 $\alpha(t)$를 다음과 같이 정의한다.
\begin{equation}
\alpha(t) = \frac{1}{2}\left(1 - \cos\frac{\pi t}{T}\right)
\label{eq:cosine_blending}
\end{equation}
이 함수는 다음 성질을 가진다.
\begin{itemize}[nosep]
  \item $\alpha(0) = 0$: 초기 시점에서 출발 궤도만 반영
  \item $\alpha(T) = 1$: 종단 시점에서 도착 궤도만 반영
  \item $\alpha'(0) = \alpha'(T) = 0$: 양 끝에서 기울기 0, 즉 경계에서 상태 변화가 부드러움
  \item $\alpha(T/2) = 1/2$: 중간 시점에서 두 궤적의 동일 가중 혼합
  \item $C^{\infty}$ 연속
\end{itemize}

선형 혼합 $\alpha_{\text{lin}}(t) = t/T$ 대비 cosine blending의 장점은 양단에서 기울기가 0이므로,
경계조건 부근의 상태 불연속이 완화된다는 점이다. 혼합된 초기값은
\[
\mathbf{x}_k^{\text{guess}} = (1 - \alpha_k)\,\mathbf{x}_k^{\text{dep}} + \alpha_k\,\mathbf{x}_k^{\text{arr}}
\]
이다. 여기서 $\mathbf{x}_k = [\mathbf{r}_k^{\top}, \mathbf{v}_k^{\top}]^{\top}$는 6차원 상태벡터이다.

\textbf{초기 제어 입력.}
초기 추력 벡터는 $\mathbf{u}_k^{\text{guess}} = \mathbf{0}$으로 설정한다.
이는 케플러 궤적(무추력 궤적)에서 출발한다는 전제와 일치하며, 솔버가 추력 프로파일을 자유롭게 탐색하도록 허용한다.
비영(nonzero) 추력 초기값은 NLP의 해공간 탐색을 특정 방향으로 편향시킬 수 있다~\cite{Betts2010}.


%% -------------------------------------------------------------------------
\subsection{케플러 전파법 (Keplerian Guess)}
\label{sec:keplerian_guess}
%% -------------------------------------------------------------------------

도착 궤도 정보를 사용하지 않고, 출발 궤도에서의 순수 케플러 전파만으로 초기값을 생성하는 단순한 방법이다.
출발 궤도요소로부터 초기 상태벡터를 구한 후, 모든 시간점에서 전방 전파한다.
\[
(\mathbf{r}_k^{\text{guess}}, \mathbf{v}_k^{\text{guess}}) = \text{Kepler}(\mathbf{r}_0, \mathbf{v}_0, t_k), \quad k = 0, 1, \ldots, N-1
\]

이 방법은 도착 궤도 경계조건을 전혀 반영하지 않으므로, 장반경 변화 $\Delta a$나 경사각 변화 $\Delta i$가 큰 문제에서는
선형 보간법보다 초기 구속조건 위반이 크다.
그러나 구현이 단순하고 항상 물리적으로 타당한 궤적을 생성한다는 장점이 있어,
선형 보간법의 대안 또는 첫 번째 시도로 사용된다.


%% -------------------------------------------------------------------------
\subsection{True Anomaly 초기값 휴리스틱}
\label{sec:nu_heuristic}
%% -------------------------------------------------------------------------

원궤도($e = 0$)에서는 근점 인수(argument of periapsis) $\omega$가 정의되지 않으므로,
true anomaly $\nu$를 argument of latitude로 취급한다.
출발 true anomaly $\nu_0$는 자유 변수이며, 그 초기값은 기동 유형에 따라 다음과 같이 설정한다.

\textbf{Coplanar 전이 ($\Delta i \approx 0$): $\nu_0 = 0$.}
경사각 변화가 없는 궤도전이에서는 장반경만 변화시키면 되므로, 추력의 주성분이 속도 방향(접선 방향)이다.
에너지 변화율은
\[
\dot{\varepsilon} = \mathbf{v} \cdot \mathbf{a}_{\text{thrust}}
\]
이므로, 속력이 큰 위치에서 접선 추력의 에너지 변환 효율이 높다 (Oberth 효과~\cite{Oberth1972}).
원궤도에서 속력은 균일하나, 전이 과정 중 생성되는 타원 궤도에서는 근지점($\nu = 0$ 부근) 속력이 최대이다.
따라서 $\nu_0 = 0$에서 출발하면 기동 초기의 추력 효율이 최적화된다.

\textbf{경사각 변경 전이 ($\Delta i \ne 0$): $\nu_0 = \pi/2$.}
경사각 변경은 궤도면 법선 방향의 속도 성분을 바꾸는 기동이다. 경사각 변경에 필요한 $\Delta v$는
\[
\Delta v_i = 2v \sin\frac{\Delta i}{2}
\]
이며, 이 기동은 line of nodes (승교점--강교점 연결선) 방향에서 가장 효율적이다~\cite{Curtis2020}.
원궤도에서 $\Omega = 0$이면 line of nodes는 $x$축 방향이고,
이에 수직인 궤도면 내 방향은 argument of latitude $u = \pi/2$이다.
보다 정확히는, line of nodes 위의 점($u = 0$ 또는 $u = \pi$)에서 면외 기동을 수행해야 하나,
연속 추력 문제에서는 기동이 분산되므로, $\nu_0 = \pi/2$를 설정하여 기동 구간의 중심이 노드 근처에 오도록 배치한다.

이 휴리스틱은 NLP의 true anomaly 자유 변수에 대한 초기값만 제공하며, 최적화 과정에서 솔버가 자유롭게 조정한다.


%% -------------------------------------------------------------------------
\subsection{Warm Start: Pass~1에서 Pass~2로의 보간}
\label{sec:warm_start}
%% -------------------------------------------------------------------------

Two-Pass 최적화에서 Pass~1 (Hermite-Simpson collocation, 균일 격자)의 해를 Pass~2 (LGL pseudospectral, 비균일 격자)의 초기값으로 사용한다.
균일 시간점에서 구한 상태/제어 변수를 LGL 노드 위치에서 cubic spline으로 보간한다.
이 과정의 상세는 B3 보고서에서 다루며, 여기서는 전략적 의의만 기술한다.

Warm start는 Pass~1에서 이미 동역학 구속조건을 근사적으로 만족하는 해를 제공하므로,
Pass~2의 초기 구속조건 위반이 크게 줄어든다.
이는 IPOPT의 실행가능성 복원 단계를 단축하고, Newton 반복 횟수를 감소시킨다~\cite{Wachter2006}.


%% =========================================================================
\section{구현 매핑}
\label{sec:implementation}
%% =========================================================================

%% -------------------------------------------------------------------------
\subsection{파일 구조}
\label{sec:file_structure}
%% -------------------------------------------------------------------------

\begin{table}[htbp]
\centering
\caption{초기값 생성 관련 파일 구조}
\label{tab:file_structure}
\begin{tabular}{@{}ll@{}}
\toprule
파일 & 역할 \\
\midrule
\texttt{optimizer/initial\_guess.py} & 초기값 생성 함수 2종 \\
\texttt{astrodynamics/kepler.py} & Universal Variable 케플러 전파 \\
\texttt{astrodynamics/orbital\_elements.py} & 궤도요소 $\leftrightarrow$ 상태벡터 변환 \\
\texttt{constants.py} & $\mu_\oplus = 398600.4418$\,km$^3$/s$^2$ \\
\texttt{types.py} & \texttt{TransferConfig} 데이터 클래스 \\
\bottomrule
\end{tabular}
\end{table}


%% -------------------------------------------------------------------------
\subsection{\texttt{linear\_interpolation\_guess} (11--79행)}
\label{sec:impl_linear}
%% -------------------------------------------------------------------------

\begin{table}[htbp]
\centering
\caption{\texttt{linear\_interpolation\_guess} 수학--코드 매핑}
\label{tab:linear_interp_mapping}
\small
\begin{tabular}{@{}llcl@{}}
\toprule
수학 표현 & Python 구현 & 행 & 비고 \\
\midrule
$\nu_0 = 0,\; \nu_f = \pi$ & \texttt{nu0\_guess = 0.0; nuf\_guess = np.pi} & 39--40 & coplanar 기준 \\
$(a_0, 0, i_0, 0, 0, \nu_0) \to (\mathbf{r}_0, \mathbf{v}_0)$ & \texttt{oe\_to\_rv(oe0, MU\_EARTH)} & 43--44 & 출발 궤도 \\
$(a_f, 0, i_f, 0, 0, \nu_f) \to (\mathbf{r}_f, \mathbf{v}_f)$ & \texttt{oe\_to\_rv(oef, MU\_EARTH)} & 47--48 & 도착 궤도 \\
$\text{Kepler}(\mathbf{r}_0, \mathbf{v}_0, t_k)$ & \texttt{kepler\_propagate(r0, v0, dt\_fwd, ...)} & 59--63 & $k=1,\ldots,N{-}1$ \\
$\text{Kepler}(\mathbf{r}_f, \mathbf{v}_f, t_k{-}T)$ & \texttt{kepler\_propagate(rf, vf, dt\_bwd, ...)} & 65--69 & $k=N{-}2,\ldots,0$ \\
$\alpha_k = \tfrac{1}{2}(1 - \cos(\pi t_k/T))$ & \texttt{0.5*(1.0-np.cos(np.pi*t[k]/T))} & 74 & cosine blending \\
$(1{-}\alpha_k)\mathbf{x}_k^{\text{dep}} + \alpha_k \mathbf{x}_k^{\text{arr}}$ & \texttt{(1-alpha)*x\_depart + alpha*x\_arrive} & 75 & 상태 혼합 \\
$\mathbf{u}_k = \mathbf{0}$ & \texttt{u\_guess = np.zeros((3, N\_points))} & 77 & 무추력 초기값 \\
\bottomrule
\end{tabular}
\end{table}

\textbf{반환값}: \texttt{(t, x\_guess, u\_guess, nu0\_guess, nuf\_guess)}.
시간 배열 \texttt{t}는 \texttt{np.linspace(0, T, N\_points)}로 균일 격자이다.


%% -------------------------------------------------------------------------
\subsection{\texttt{keplerian\_guess} (82--125행)}
\label{sec:impl_keplerian}
%% -------------------------------------------------------------------------

\begin{table}[htbp]
\centering
\caption{\texttt{keplerian\_guess} 수학--코드 매핑}
\label{tab:keplerian_mapping}
\small
\begin{tabular}{@{}llcl@{}}
\toprule
수학 표현 & Python 구현 & 행 & 비고 \\
\midrule
경사각 변경 판별 & \texttt{abs(config.delta\_i) > 1e-10} & 104 & 수치 허용값 \\
$\nu_0 = \pi/2$ 또는 $0$ & \texttt{np.pi/2 if is\_plane\_change else 0.0} & 105 & \S\ref{sec:nu_heuristic} 휴리스틱 \\
$(a_0, 0, i_0, 0, 0, \nu_0) \to (\mathbf{r}_0, \mathbf{v}_0)$ & \texttt{oe\_to\_rv(oe0, MU\_EARTH)} & 109--110 & 출발 궤도 \\
$\text{Kepler}(\mathbf{r}_0, \mathbf{v}_0, t_k)$ & \texttt{kepler\_propagate(r0, v0, dt, ...)} & 117--121 & $k=1,\ldots,N{-}1$ \\
$\mathbf{u}_k = \mathbf{0}$ & \texttt{u\_guess = np.zeros((3, N\_points))} & 123 & 무추력 초기값 \\
\bottomrule
\end{tabular}
\end{table}

이 함수는 도착 궤도 정보를 사용하지 않으므로 $\nu_f$는 $\pi$로 고정 반환한다 (106행).


%% -------------------------------------------------------------------------
\subsection{\texttt{kepler\_propagate} (\texttt{kepler.py}, 31--144행)}
\label{sec:impl_kepler}
%% -------------------------------------------------------------------------

\begin{table}[htbp]
\centering
\caption{\texttt{kepler\_propagate} 수학--코드 매핑}
\label{tab:kepler_propagate_mapping}
\small
\begin{tabular}{@{}llc@{}}
\toprule
수학 표현 & Python 구현 & 행 \\
\midrule
$\alpha = 2/r_0 - v_0^2/\mu$ & \texttt{2.0/r0\_mag - v0\_mag**2/mu} & 63 \\
$c_2(\psi)$ & \texttt{\_stumpff\_c2(psi)} & 9--16 \\
$c_3(\psi)$ & \texttt{\_stumpff\_c3(psi)} & 19--28 \\
$F(\chi) = 0$ & \texttt{chi3*c3 + rdotv*chi2*c2 + ...} & 103--105 \\
$F'(\chi)$ & \texttt{chi2*c2 + rdotv*chi*(1-psi*c3) + ...} & 108--110 \\
$\delta\chi = -F/F'$ & \texttt{delta\_chi = -fn / r\_chi} & 115 \\
$f = 1 - \chi^2 c_2 / r_0$ & \texttt{1.0 - chi2/r0\_mag*c2} & 130 \\
$g = \Delta t - \chi^3 c_3 / \sqrt{\mu}$ & \texttt{dt - chi3/sqrt\_mu*c3} & 131 \\
$\dot{f} = \sqrt{\mu}\,\chi(\psi c_3 - 1)/(r\,r_0)$ & \texttt{sqrt\_mu/(r\_mag*r0\_mag)*chi*(psi*c3-1)} & 138 \\
$\dot{g} = 1 - \chi^2 c_2 / r$ & \texttt{1.0 - chi2/r\_mag*c2} & 139 \\
$\mathbf{r} = f\mathbf{r}_0 + g\mathbf{v}_0$ & \texttt{f*r0 + g*v0} & 134 \\
$\mathbf{v} = \dot{f}\mathbf{r}_0 + \dot{g}\mathbf{v}_0$ & \texttt{fdot*r0 + gdot*v0} & 142 \\
\bottomrule
\end{tabular}
\end{table}

Newton-Raphson 반복은 최대 100회, 허용 오차 $\varepsilon = 10^{-12}$로 설정되어 있다 (87--88행).


%% -------------------------------------------------------------------------
\subsection{설계 결정}
\label{sec:design_decisions}
%% -------------------------------------------------------------------------

\begin{enumerate}[nosep]
  \item \textbf{Cosine blending 선택}: 단순 선형 보간($\alpha = t/T$) 대비 양 끝에서 기울기가 0이므로, 경계조건 부근의 상태 변화가 부드럽다. 이는 NLP 솔버의 초기 Jacobian 평가에서 수치적 안정성을 개선한다.

  \item \textbf{역전파 사용}: 도착 궤도에서 음의 시간($t_k - T < 0$)으로 케플러 전파하여 역방향 궤적을 생성한다. Universal Variable Formulation은 $\Delta t < 0$을 자연스럽게 처리하므로 별도의 역전파 알고리즘이 불필요하다.

  \item \textbf{True anomaly 자유 변수}: $\nu_0$, $\nu_f$는 NLP의 자유 변수이다. 초기값 함수는 이들의 초기 추정값만 반환하며, 최종 최적값은 솔버가 결정한다.

  \item \textbf{초기 추력 0 설정}: 무추력 케플러 궤적에서 출발하므로 $\mathbf{u}_k = \mathbf{0}$이 물리적으로 일관된다. 또한 비용함수 $J = \int \|\mathbf{u}\|^2\,dt$의 관점에서 초기 비용이 0이 되어, 솔버가 비용을 증가시키는 방향으로만 추력을 추가하게 된다.
\end{enumerate}


%% =========================================================================
\section{수치 검증}
\label{sec:verification}
%% =========================================================================

%% -------------------------------------------------------------------------
\subsection{경계 궤도 반경 일치}
\label{sec:verify_boundary}
%% -------------------------------------------------------------------------

선형 보간법의 초기값이 경계조건을 올바르게 반영하는지 검증한다.

\textbf{출발 경계 ($k = 0$).} $\alpha(0) = 0$이므로
\[
\mathbf{x}_0^{\text{guess}} = \mathbf{x}_0^{\text{dep}}
\]
이며, 이는 출발 궤도요소에서 변환한 상태벡터와 정확히 일치한다. 따라서
\[
\|\mathbf{r}_0^{\text{guess}}\| = a_0 \cdot \frac{1 - e_0^2}{1 + e_0 \cos\nu_0}
\]
이고, 원궤도($e_0 = 0$)에서 $\|\mathbf{r}_0\| = a_0 = R_\oplus + h_0$이다.

\textbf{도착 경계 ($k = N-1$).} $\alpha(T) = 1$이므로
\[
\mathbf{x}_{N-1}^{\text{guess}} = \mathbf{x}_{N-1}^{\text{arr}}
\]
이며, $\|\mathbf{r}_{N-1}\| = a_f = a_0 + \Delta a$이다.

\textbf{검증 기준}: $\left| \|\mathbf{r}_0\| - a_0 \right| < 10^{-10}$\,km, $\left| \|\mathbf{r}_{N-1}\| - a_f \right| < 10^{-10}$\,km.


%% -------------------------------------------------------------------------
\subsection{True Anomaly 물리적 타당성}
\label{sec:verify_nu}
%% -------------------------------------------------------------------------

\begin{table}[htbp]
\centering
\caption{True anomaly 초기값 검증}
\label{tab:nu_verification}
\begin{tabular}{@{}llll@{}}
\toprule
기동 유형 & $\nu_0$ 초기값 & 물리적 근거 & 검증 기준 \\
\midrule
Coplanar ($\Delta i = 0$) & $0$ & Oberth 효과: 근지점에서 접선 기동 효율 최대 & $\nu_0 = 0.0$ 확인 \\
Plane change ($\Delta i \ne 0$) & $\pi/2$ & Line of nodes에서 면외 기동 효율 최대 & $\nu_0 = \pi/2$ 확인 \\
선형 보간법 (항상) & $\nu_f = \pi$ & 반궤도 전이: 대칭적 기동 분포 & $\nu_f = \pi$ 확인 \\
\bottomrule
\end{tabular}
\end{table}


%% -------------------------------------------------------------------------
\subsection{Cosine Blending 함수 성질 검증}
\label{sec:verify_blending}
%% -------------------------------------------------------------------------

$\alpha(t) = \frac{1}{2}(1 - \cos(\pi t/T))$에 대해 다음을 수치적으로 확인한다.

\begin{table}[htbp]
\centering
\caption{Cosine blending 함수 성질 검증}
\label{tab:blending_verification}
\begin{tabular}{@{}lll@{}}
\toprule
성질 & 기대값 & 허용 오차 \\
\midrule
$\alpha(0)$ & $0$ & $< 10^{-15}$ \\
$\alpha(T)$ & $1$ & $< 10^{-15}$ \\
$\alpha(T/2)$ & $0.5$ & $< 10^{-15}$ \\
$\alpha'(0) = \frac{\pi}{2T}\sin(0)$ & $0$ & $< 10^{-15}$ \\
$\alpha'(T) = \frac{\pi}{2T}\sin(\pi)$ & $0$ & $< 10^{-15}$ \\
$\alpha(t) \in [0, 1]$ for all $t \in [0, T]$ & 단조 증가 & 모든 격자점 확인 \\
\bottomrule
\end{tabular}
\end{table}


%% -------------------------------------------------------------------------
\subsection{케플러 전파 보존량 검증}
\label{sec:verify_kepler}
%% -------------------------------------------------------------------------

Universal Variable Formulation에 의한 케플러 전파가 이체 문제의 보존량을 유지하는지 검증한다.

\textbf{에너지 보존.}
비역학 에너지 $\varepsilon = v^2/2 - \mu/r$은 궤도상 어디서나 일정하다.
\[
\left| \varepsilon(t_k) - \varepsilon(0) \right| < 10^{-10} \text{\,km}^2/\text{s}^2, \quad \forall k
\]

\textbf{각운동량 보존.}
비각운동량 $\mathbf{h} = \mathbf{r} \times \mathbf{v}$의 크기는 일정하다.
\[
\left| \|\mathbf{h}(t_k)\| - \|\mathbf{h}(0)\| \right| < 10^{-10} \text{\,km}^2/\text{s}, \quad \forall k
\]

\textbf{Lagrange 계수 항등식.}
$f\dot{g} - \dot{f}g = 1$이 수치적으로 만족되는지 확인한다.


%% -------------------------------------------------------------------------
\subsection{검증 결과 요약}
\label{sec:verify_summary}
%% -------------------------------------------------------------------------

\begin{table}[htbp]
\centering
\caption{수치 검증 결과 요약}
\label{tab:verification_summary}
\begin{tabular}{@{}llll@{}}
\toprule
검증 항목 & 기대값 & 허용 오차 & 결과 \\
\midrule
출발점 반경 $\|\mathbf{r}_0\| = a_0$ & $R_\oplus + h_0$ & $10^{-10}$\,km & 통과 \\
도착점 반경 $\|\mathbf{r}_{N-1}\| = a_f$ & $a_0 + \Delta a$ & $10^{-10}$\,km & 통과 \\
$\nu_0$ 휴리스틱 (coplanar) & $0.0$\,rad & 정확 일치 & 통과 \\
$\nu_0$ 휴리스틱 (plane change) & $\pi/2$\,rad & 정확 일치 & 통과 \\
Blending $\alpha(0) = 0$, $\alpha(T) = 1$ & 경계값 & $10^{-15}$ & 통과 \\
케플러 전파 에너지 보존 & $\Delta\varepsilon = 0$ & $10^{-10}$\,km$^2$/s$^2$ & 통과 \\
케플러 전파 각운동량 보존 & $\Delta h = 0$ & $10^{-10}$\,km$^2$/s & 통과 \\
\bottomrule
\end{tabular}
\end{table}


%% =========================================================================
\bibliographystyle{IEEEtran}
\bibliography{references}

\end{document}
